\section{Алгебры с одной операцией. Группы. Группа перестановок.}

\subsection*{Определение группы}

\paragraph{Определение}
\textbf{Группа} — это моноид, в котором для каждого элемента $a$ существует обратный элемент $a^{-1}$ такой, что:
$$\forall a \in G \ (\exists a^{-1} \in G: a \ast a^{-1} = a^{-1} \ast a = e).$$
Здесь $e$ — единица моноида, а операция $\ast$ часто называется умножением.

\paragraph{Основные понятия}
\begin{itemize}
    \item \textbf{Порядок группы:} Мощность носителя $|G|$.
    \item \textbf{Абелева группа:} Группа, в которой операция коммутативна ($\forall a, b: a \ast b = b \ast a$). В таких группах часто используют аддитивную запись: операция $+$, нейтральный элемент $0$, обратный $-a$.
\end{itemize}

\subsection*{Три теоремы о свойствах групп}

\paragraph{Теорема 1 (Единственность обратного)}
В группе обратный элемент $a^{-1}$ для любого элемента $a$ единственен.

Пусть существуют два обратных элемента $a^{-1}$ и $b$ для элемента $a$. Тогда $a \ast a^{-1} = a^{-1} \ast a = e$ и $a \ast b = b \ast a = e$.
Имеем: $a^{-1} = a^{-1} \ast e = a^{-1} \ast (a \ast b) = (a^{-1} \ast a) \ast b = e \ast b = b$. Следовательно, $a^{-1} = b$.

\paragraph{Теорема 2 (Соотношения в группе)}
В любой группе выполняются следующие свойства:
\begin{enumerate}
    \item $(a \ast b)^{-1} = b^{-1} \ast a^{-1}$.
    \item Закон сокращения слева: $a \ast b = a \ast c \implies b = c$.
    \item Закон сокращения справа: $b \ast a = c \ast a \implies b = c$.
    \item Двойной обратный: $(a^{-1})^{-1} = a$.
\end{enumerate}

1. Проверим $(a \ast b) \ast (b^{-1} \ast a^{-1}) = a \ast (b \ast b^{-1}) \ast a^{-1} = a \ast e \ast a^{-1} = a \ast a^{-1} = e$. Аналогично для умножения слева.
2. Если $a \ast b = a \ast c$, то умножим обе части слева на $a^{-1}$: $a^{-1} \ast (a \ast b) = a^{-1} \ast (a \ast c) \implies (a^{-1} \ast a) \ast b = (a^{-1} \ast a) \ast c \implies e \ast b = e \ast c \implies b = c$.
3. Аналогично пункту 2, умножением на $a^{-1}$ справа.
4. По определению $a^{-1} \ast a = e$, значит $a$ является обратным для $a^{-1}$. Так как обратный единственен, $(a^{-1})^{-1} = a$.

\paragraph{Теорема 3 (Разрешимость уравнений)}
В группе можно однозначно решить уравнение $a \ast x = b$.

Решение существует: подставим $x = a^{-1} \ast b$. Тогда $a \ast (a^{-1} \ast b) = (a \ast a^{-1}) \ast b = e \ast b = b$.
Решение единственно: если $x_1$ и $x_2$ — решения, то $a \ast x_1 = b$ и $a \ast x_2 = b$. Значит $a \ast x_1 = a \ast x_2$, и по закону сокращения $x_1 = x_2$.

\subsection*{Группа перестановок}

\paragraph{Определение}
Биективная функция $f: X \to X$ называется \textbf{перестановкой}. На конечном множестве перестановку удобно задавать таблицей из двух строк, называемой \textbf{подстановкой}:
$$\sigma = \begin{pmatrix} 1 & 2 & \dots & n \\ f(1) & f(2) & \dots & f(n) \end{pmatrix}$$
Множество всех перестановок $n$-элементного множества образует \textbf{симметрическую группу} $S_n$.

\paragraph{Свойства $S_n$}
\begin{itemize}
    \item \textbf{Операция:} Суперпозиция функций (результат применения $f$, а затем $g$ обозначается $fg$).
    \item \textbf{Единица:} Тождественная перестановка $e(x) = x$.
    \item \textbf{Обратный элемент:} Обратная функция $f^{-1}$, таблицу которой можно получить, поменяв строки исходной таблицы местами (и отсортировав верхнюю).
    \item \textbf{Порядок:} $|S_n| = n!$.
\end{itemize}

\subsection*{Примеры групп}
\begin{enumerate}
    \item $(\mathbb{Z}; +)$ — абелева группа целых чисел. $e=0, a^{-1}=-a$.
    \item $(\mathbb{Q}^+; \cdot)$ — абелева группа положительных рациональных чисел. $e=1, (m/n)^{-1} = n/m$.
    \item $(2^M; \Delta)$ — булеан с операцией симметрической разности. $e=\varnothing, a^{-1}=a$.
    \item Группа невырожденных матриц $n \times n$ относительно умножения.
\end{enumerate}